% !TEX program = xelatex
% ---------------------------------------------------------------------------
% VĂN PHẠM PHI NGỮ CẢNH CỦA NGÔN NGỮ UPL
% Biên dịch: XeLaTeX (Overleaf: Menu -> Compiler -> XeLaTeX; dòng magic comment
% ở trên đã tự chọn XeLaTeX). Cần biên dịch 2 lần để mục lục/tham chiếu đúng.
% ---------------------------------------------------------------------------
\documentclass[11pt,a4paper]{article}

% ---------- Phông chữ (hỗ trợ đầy đủ tiếng Việt) ----------
\usepackage{fontspec}
\setmainfont{TeX Gyre Pagella}
\setsansfont{TeX Gyre Heros}
\setmonofont{Latin Modern Mono}[Scale=MatchLowercase]

% ---------- Trang giấy, tiêu đề ----------
\usepackage[a4paper,margin=2.5cm,headheight=22pt]{geometry}
\usepackage{microtype}
\usepackage{xcolor}
\usepackage{titlesec}
\usepackage{fancyhdr}
\usepackage{enumitem}
\usepackage[hidelinks]{hyperref}

% ---------- Bảng, công thức, mã ----------
\usepackage{array}
\usepackage{booktabs}
\usepackage{tabularx}
\usepackage{longtable}
\usepackage{amsmath}
\usepackage{amssymb}
\usepackage{listings}
\usepackage[edges]{forest}

% ---------- Màu và kiểu chữ ----------
\definecolor{uplblue}{HTML}{1F3A5F}
\definecolor{uplgray}{HTML}{F3F5F8}
\definecolor{uplrule}{HTML}{C9D1DC}
\definecolor{upltok}{HTML}{8A2D1F}

\titleformat{\section}
  {\sffamily\Large\bfseries\color{uplblue}}
  {\thesection.}{0.6em}{}[\vspace{-0.4em}{\color{uplrule}\titlerule}]
\titlespacing*{\section}{0pt}{2.0ex plus 0.5ex}{1.2ex}

\pagestyle{fancy}
\fancyhf{}
\renewcommand{\headrulewidth}{0.3pt}
\renewcommand{\headrule}{\hbox to\headwidth{\color{uplrule}\leaders\hrule height \headrulewidth\hfill}}
\lhead{\small\sffamily\color{uplblue} Văn phạm phi ngữ cảnh của ngôn ngữ UPL}
\rhead{\small\sffamily\color{uplblue} BTL01}
\cfoot{\small\thepage}

\setlength{\parindent}{0pt}
\setlength{\parskip}{0.55em}
\renewcommand{\arraystretch}{1.25}
\setlist[itemize]{itemsep=0.25em,topsep=0.3em}
\setlist[enumerate]{itemsep=0.25em,topsep=0.3em}

% ---------- Macro cho ký hiệu văn phạm ----------
% \nt{X}: ký hiệu không kết thúc;  \tm{x}: ký hiệu kết thúc (token)
\newcommand{\nt}[1]{\textit{#1}}
\newcommand{\tm}[1]{\texttt{\textcolor{upltok}{#1}}}
\newcommand{\eps}{\ensuremath{\varepsilon}}
\newcommand{\ar}{\ensuremath{\rightarrow}}
\newcommand{\alt}{\ensuremath{\mid}}
\newcommand{\dv}{\ensuremath{\Rightarrow}}
% \rn{k}: số hiệu luật sinh, ví dụ (12)
\newcommand{\rn}[1]{(#1)}

% ---------- Mã nguồn UPL ----------
\lstdefinelanguage{UPL}{
  morekeywords={begin,end,int,bool,if,then,else,do,while,for,print,true,false},
  sensitive=true,
  morecomment=[l]{//},
  morecomment=[s]{/*}{*/},
}
\lstset{
  language=UPL,
  basicstyle=\ttfamily\small,
  keywordstyle=\color{uplblue},
  commentstyle=\itshape\color{gray},
  backgroundcolor=\color{uplgray},
  frame=single,
  rulecolor=\color{uplrule},
  framesep=6pt,
  xleftmargin=0pt,
  columns=fullflexible,
  keepspaces=true,
  showstringspaces=false,
}

\begin{document}

% ======================= Tiêu đề =======================
\begin{flushleft}
  {\sffamily\bfseries\fontsize{22}{27}\selectfont\color{uplblue}
   VĂN PHẠM PHI NGỮ CẢNH CỦA\\[2pt] NGÔN NGỮ UPL\par}
  \vspace{0.8em}
  {\color{uplrule}\rule{\linewidth}{0.8pt}}\par
  \vspace{0.2em}
  \begin{tabular}{@{}l@{\hspace{0.6em}}l@{}}
    24021485 & Lã Việt Hoàng \\
    24021533 & Nguyễn Đức Khiêm \\
  \end{tabular}
\end{flushleft}

% ======================= 1. Giới thiệu =======================
\section{Giới thiệu}

Tài liệu này đặc tả văn phạm phi ngữ cảnh (CFG) của ngôn ngữ lập trình UPL.
Văn phạm là bộ bốn $G = (N, T, P, S)$, trong đó $N$ là tập ký hiệu không kết
thúc, $T$ là tập ký hiệu kết thúc (token), $P$ là tập luật sinh và
$S = \nt{Program}$ là ký hiệu bắt đầu.

% ======================= 2. Giả định thiết kế =======================
\section{Các giả định thiết kế}

Đề bài chưa quy định một số chi tiết, nên văn phạm này thống nhất như sau:

\begin{enumerate}
  \item Câu lệnh \tm{for} có dạng kiểu C:
        \texttt{for (khởi\_tạo; điều\_kiện; cập\_nhật) \{ \dots{} \}}.
        Phần khởi tạo là khai báo hoặc phép gán, phần cập nhật là phép gán.
  \item Bổ sung hai hằng boolean \tm{true} và \tm{false}.
  \item Thân của \tm{if}, \tm{else}, \tm{do-while}, \tm{for} bắt buộc nằm
        trong cặp \texttt{\{ \}}.
  \item Phép so sánh (\tm{>}, \tm{>=}, \tm{==}) không có tính kết hợp, nghĩa
        là \texttt{a > b > c} không hợp lệ.
  \item Văn phạm chỉ mô tả cú pháp. Các ràng buộc ngữ nghĩa (khai báo trước
        khi dùng, khai báo lại, tương thích kiểu) không thuộc phạm vi tài
        liệu này.
\end{enumerate}

% ======================= 3. Token =======================
\section{Ký hiệu kết thúc (token)}

\begin{table}[h]
\centering
\begin{tabularx}{\linewidth}{@{}l >{\raggedright\arraybackslash}p{4.6cm} X@{}}
\toprule
\textbf{Nhóm} & \textbf{Token} & \textbf{Mô tả} \\
\midrule
Từ khóa & \tm{begin} \tm{end} \tm{int} \tm{bool} \tm{if} \tm{then} \tm{else}
          \tm{do} \tm{while} \tm{for} \tm{print} \tm{true} \tm{false}
        & Từ khóa dành riêng, không dùng làm định danh \\
\addlinespace
Định danh & \tm{id}
          & Bắt đầu bằng chữ cái, nếu có chữ số thì chỉ nằm ở cuối:
            \texttt{[A-Za-z]+[0-9]*} \\
\addlinespace
Hằng số nguyên & \tm{num}
               & Dãy một hoặc nhiều chữ số: \texttt{[0-9]+} \\
\addlinespace
Toán tử số học & \tm{+} \tm{*} & Cộng, nhân \\
\addlinespace
Toán tử so sánh & \tm{>} \tm{>=} \tm{==}
                & Lớn hơn, lớn hơn hoặc bằng, bằng \\
\addlinespace
Phép gán & \tm{=} & Gán giá trị \\
\addlinespace
Phân cách & \tm{(} \tm{)} \tm{\{} \tm{\}} \tm{;} & \\
\bottomrule
\end{tabularx}
\end{table}

Khoảng trắng và comment kiểu Java (\texttt{//} \dots{} đến cuối dòng,
\texttt{/*} \dots{} \texttt{*/} có thể nhiều dòng) được bộ phân tích từ vựng
loại bỏ, không xuất hiện trong văn phạm.

% ======================= 4. Không kết thúc =======================
\section{Ký hiệu không kết thúc}

\begin{table}[h]
\centering
\begin{tabularx}{\linewidth}{@{}>{\itshape}l X@{}}
\toprule
\textbf{Ký hiệu} & \textbf{Ý nghĩa} \\
\midrule
Program  & Toàn bộ chương trình (ký hiệu bắt đầu) \\
StmtList & Danh sách câu lệnh (có thể rỗng) \\
Stmt     & Một câu lệnh \\
Block    & Khối lệnh trong \texttt{\{ \}} \\
Decl     & Khai báo biến (có thể kèm khởi tạo) \\
Type     & Kiểu dữ liệu \\
Assign   & Câu lệnh gán \\
Print    & Câu lệnh in \\
IfStmt   & Câu lệnh rẽ nhánh \\
DoWhile  & Câu lệnh lặp do-while \\
ForStmt  & Câu lệnh lặp for \\
ForInit  & Phần khởi tạo của for \\
Expr     & Biểu thức (có thể có so sánh) \\
RelOp    & Toán tử so sánh \\
Arith    & Biểu thức cộng \\
Term     & Biểu thức nhân \\
Factor   & Toán hạng cơ bản \\
\bottomrule
\end{tabularx}
\end{table}

% ======================= 5. Luật sinh =======================
\clearpage % giữ trọn 37 luật sinh trên cùng một trang
\section{Các luật sinh}

\begin{longtable}{@{}r@{\hspace{0.9em}}l@{\hspace{0.7em}}c@{\hspace{0.7em}}l@{}}
% --- Chương trình và danh sách câu lệnh ---
\rn{1}  & \nt{Program}  & \ar  & \tm{begin} \nt{StmtList} \tm{end} \\
\addlinespace
\rn{2}  & \nt{StmtList} & \ar  & \nt{StmtList} \nt{Stmt} \\
\rn{3}  &               & \alt & \eps \\
\addlinespace
% --- Câu lệnh ---
\rn{4}  & \nt{Stmt}     & \ar  & \nt{Decl} \tm{;} \\
\rn{5}  &               & \alt & \nt{Assign} \tm{;} \\
\rn{6}  &               & \alt & \nt{Print} \tm{;} \\
\rn{7}  &               & \alt & \nt{IfStmt} \\
\rn{8}  &               & \alt & \nt{DoWhile} \\
\rn{9}  &               & \alt & \nt{ForStmt} \\
\rn{10} &               & \alt & \nt{Block} \\
\addlinespace
\rn{11} & \nt{Block}    & \ar  & \tm{\{} \nt{StmtList} \tm{\}} \\
\addlinespace
% --- Khai báo, gán, in ---
\rn{12} & \nt{Decl}     & \ar  & \nt{Type} \tm{id} \\
\rn{13} &               & \alt & \nt{Type} \tm{id} \tm{=} \nt{Expr} \\
\rn{14} & \nt{Type}     & \ar  & \tm{int} \\
\rn{15} &               & \alt & \tm{bool} \\
\rn{16} & \nt{Assign}   & \ar  & \tm{id} \tm{=} \nt{Expr} \\
\rn{17} & \nt{Print}    & \ar  & \tm{print} \tm{(} \nt{Expr} \tm{)} \\
\addlinespace
% --- Rẽ nhánh và vòng lặp ---
\rn{18} & \nt{IfStmt}   & \ar  & \tm{if} \tm{(} \nt{Expr} \tm{)} \tm{then} \nt{Block} \\
\rn{19} &               & \alt & \tm{if} \tm{(} \nt{Expr} \tm{)} \tm{then} \nt{Block} \tm{else} \nt{Block} \\
\rn{20} & \nt{DoWhile}  & \ar  & \tm{do} \nt{Block} \tm{while} \tm{(} \nt{Expr} \tm{)} \tm{;} \\
\rn{21} & \nt{ForStmt}  & \ar  & \tm{for} \tm{(} \nt{ForInit} \tm{;} \nt{Expr} \tm{;} \nt{Assign} \tm{)} \nt{Block} \\
\rn{22} & \nt{ForInit}  & \ar  & \nt{Decl} \\
\rn{23} &               & \alt & \nt{Assign} \\
\addlinespace
% --- Biểu thức ---
\rn{24} & \nt{Expr}     & \ar  & \nt{Arith} \\
\rn{25} &               & \alt & \nt{Arith} \nt{RelOp} \nt{Arith} \\
\rn{26} & \nt{RelOp}    & \ar  & \tm{>} \\
\rn{27} &               & \alt & \tm{>=} \\
\rn{28} &               & \alt & \tm{==} \\
\addlinespace
\rn{29} & \nt{Arith}    & \ar  & \nt{Arith} \tm{+} \nt{Term} \\
\rn{30} &               & \alt & \nt{Term} \\
\rn{31} & \nt{Term}     & \ar  & \nt{Term} \tm{*} \nt{Factor} \\
\rn{32} &               & \alt & \nt{Factor} \\
\addlinespace
\rn{33} & \nt{Factor}   & \ar  & \tm{id} \\
\rn{34} &               & \alt & \tm{num} \\
\rn{35} &               & \alt & \tm{true} \\
\rn{36} &               & \alt & \tm{false} \\
\rn{37} &               & \alt & \tm{(} \nt{Expr} \tm{)} \\
\end{longtable}

% ======================= 6. Độ ưu tiên =======================
\section{Độ ưu tiên và tính kết hợp}

\begin{table}[h]
\centering
\begin{tabularx}{\linewidth}{@{}c c l X@{}}
\toprule
\textbf{Mức (cao đến thấp)} & \textbf{Toán tử} & \textbf{Kết hợp}
  & \textbf{Được thể hiện bởi} \\
\midrule
1 & \tm{(\,)}        &              & \nt{Factor} \ar{} \tm{(} \nt{Expr} \tm{)} \\
2 & \tm{*}           & Trái         & \nt{Term} \ar{} \nt{Term} \tm{*} \nt{Factor} \\
3 & \tm{+}           & Trái         & \nt{Arith} \ar{} \nt{Arith} \tm{+} \nt{Term} \\
4 & \tm{>} \tm{>=} \tm{==} & Không kết hợp
  & \nt{Expr} \ar{} \nt{Arith} \nt{RelOp} \nt{Arith} \\
\bottomrule
\end{tabularx}
\end{table}

\begin{itemize}
  \item Toán tử nằm càng sâu trong chuỗi
        $\nt{Expr} \to \nt{Arith} \to \nt{Term} \to \nt{Factor}$ thì ưu tiên
        càng cao. Vì \tm{*} nằm sâu hơn \tm{+} nên \tm{*} ưu tiên hơn \tm{+}.
  \item Đệ quy trái làm \tm{+} và \tm{*} kết hợp trái: \texttt{a+b+c} được
        hiểu là \texttt{(a+b)+c}.
  \item Luật \rn{25} chỉ chứa một \nt{RelOp} nên không thể sinh
        \texttt{a>b>c}.
  \item Dấu ngoặc ở \nt{Factor} cho phép đặt lại thứ tự tính toán.
\end{itemize}

% ======================= 7. Không nhập nhằng =======================
\section{Tính không nhập nhằng}

\begin{itemize}
  \item \textbf{Dangling else:} thân của \tm{if} luôn là \nt{Block} (bắt đầu
        bằng \texttt{\{} và kết thúc bằng \texttt{\}}), nên \tm{else} luôn
        thuộc về \tm{if} gần nhất chưa có \tm{else} mà không cần quy tắc phụ.
  \item \textbf{Biểu thức:} phân tầng \nt{Expr}/\nt{Arith}/\nt{Term}/\nt{Factor}
        cho mỗi biểu thức đúng một cây phân tích.
  \item \textbf{Câu lệnh:} mỗi dạng câu lệnh bắt đầu bằng một token khác nhau
        (\tm{int}/\tm{bool}, \tm{id}, \tm{print}, \tm{if}, \tm{do}, \tm{for},
        \texttt{\{}) nên không có hai luật cùng dẫn đến một chuỗi.
\end{itemize}

% ======================= 8. Ví dụ dẫn xuất =======================
\section{Ví dụ dẫn xuất}

\textbf{Ví dụ 1.} Chương trình rỗng \texttt{begin end}:
\[
  \nt{Program} \dv \tm{begin}\ \nt{StmtList}\ \tm{end}
               \dv \tm{begin}\ \eps\ \tm{end}
\]

\textbf{Ví dụ 2.} Câu lệnh \texttt{int y = x + 1 ;} (dẫn xuất trái nhất):
\begin{align*}
  \nt{Stmt} &\dv \nt{Decl}\ \tm{;} \\
            &\dv \nt{Type}\ \tm{id}\ \tm{=}\ \nt{Expr}\ \tm{;} \\
            &\dv \tm{int}\ \tm{id}\ \tm{=}\ \nt{Expr}\ \tm{;} \\
            &\dv \tm{int}\ \tm{id}\ \tm{=}\ \nt{Arith}\ \tm{;} \\
            &\dv \tm{int}\ \tm{id}\ \tm{=}\ \nt{Arith}\ \tm{+}\ \nt{Term}\ \tm{;} \\
            &\dv \tm{int}\ \tm{id}\ \tm{=}\ \nt{Term}\ \tm{+}\ \nt{Term}\ \tm{;} \\
            &\dv \tm{int}\ \tm{id}\ \tm{=}\ \nt{Factor}\ \tm{+}\ \nt{Term}\ \tm{;} \\
            &\dv \tm{int}\ \tm{id}\ \tm{=}\ \tm{id}\ \tm{+}\ \nt{Term}\ \tm{;} \\
            &\dv \tm{int}\ \tm{id}\ \tm{=}\ \tm{id}\ \tm{+}\ \nt{Factor}\ \tm{;} \\
            &\dv \tm{int}\ \tm{id}\ \tm{=}\ \tm{id}\ \tm{+}\ \tm{num}\ \tm{;}
\end{align*}

\textbf{Ví dụ 3.} Thứ tự ưu tiên của \texttt{x + 2 * 3}. Cây phân tích thể hiện
\texttt{2 * 3} được kết hợp trước:

\begin{center}
\begin{forest}
  for tree={
    font=\small\itshape,
    s sep=14pt, l sep=14pt,
    edge={color=uplblue},
    inner sep=2pt,
  },
  [Arith
    [Arith [Term [Factor [{\upshape\tm{id}\,(x)}]]]]
    [{\upshape\tm{+}}]
    [Term
      [Term [Factor [{\upshape\tm{num}\,(2)}]]]
      [{\upshape\tm{*}}]
      [Factor [{\upshape\tm{num}\,(3)}]]
    ]
  ]
\end{forest}
\end{center}

\textbf{Ví dụ 4.} Một chương trình có rẽ nhánh và vòng lặp:

\begin{lstlisting}
begin
  int x = 1;
  if (x >= 1) then { print(x); } else { x = x + 1; }
  do { x = x * 2; } while (x > 10);
end
\end{lstlisting}

Chương trình trên được sinh như sau:
\begin{itemize}
  \item Luật \rn{1} cho khung \texttt{begin} \dots{} \texttt{end}.
  \item Các luật \rn{2}, \rn{4}, \rn{13} cho khai báo \texttt{int x = 1;}.
  \item Các luật \rn{7}, \rn{19}, \rn{25}, \rn{27} cho
        \texttt{if} \dots{} \texttt{then} \dots{} \texttt{else} \dots{}
        với phép so sánh \texttt{>=}.
  \item Các luật \rn{8}, \rn{20}, \rn{25}, \rn{26} cho vòng
        \texttt{do} \dots{} \texttt{while}.
\end{itemize}

% ======================= 9. Đối chiếu =======================
\section{Đối chiếu với chương trình mẫu của đề}

\begin{table}[h]
\centering
\begin{tabularx}{\linewidth}{@{}>{\ttfamily\raggedright\arraybackslash}p{6.2cm} X@{}}
\toprule
\textrm{\textbf{Đoạn mã trong đề}} & \textbf{Luật áp dụng} \\
\midrule
int x;                       & \rn{4}, \rn{12}, \rn{14} \\
int y=x+1;                   & \rn{4}, \rn{13}, \rn{29}, \rn{30}, \rn{33}, \rn{34} \\
if (x>a) then\{...\}else\{...\} & \rn{7}, \rn{19}, \rn{25}, \rn{26} \\
if (x>=a) then\{...\}        & \rn{7}, \rn{18}, \rn{25}, \rn{27} \\
bool x=a==b;                 & \rn{4}, \rn{13}, \rn{15}, \rn{25}, \rn{28} \\
do\{...\}while(a>1);         & \rn{8}, \rn{20}, \rn{25}, \rn{26} \\
a=(b+10)*b;                  & \rn{5}, \rn{16}, \rn{31}, \rn{37} \\
print(a); / print(a+1);      & \rn{6}, \rn{17} \\
\bottomrule
\end{tabularx}
\end{table}

Toàn bộ chương trình mẫu của đề đều được văn phạm sinh ra. Các chỗ như
\texttt{x>a} (so sánh \tm{int} với \tm{bool}) hay khai báo lại \texttt{x} hợp
lệ về cú pháp và sẽ được phát hiện ở giai đoạn phân tích ngữ nghĩa nếu mở rộng
sau này.

\end{document}
